\documentclass[11pt]{article}

\usepackage{amsmath, amsthm, amssymb, amsfonts, amscd}
\usepackage{graphics}
\usepackage{color}

\numberwithin{equation}{section}

\theoremstyle{plain}
\newtheorem{theorem}{Theorem}[section]
\newtheorem{lemma}[theorem]{Lemma}
\newtheorem{proposition}[theorem]{Proposition}
\newtheorem{definition}[theorem]{Definition}

\newcommand{\R}{\mathbb{R}}

\bibliographystyle{plain}
\begin{document}

\title{Style Transfer in Natural Language Processing}
\author{Angelo Amato, Scott Ladenheim, Sean Salvador}
\date{\today}
\maketitle

\section{Preliminaries}

In this section, we introduce notation and concepts that will be used throughout the report. This information can all be found in standard textbooks on Probability and Statistics, see, e.g., \cite{degroot_probability_2012, wasserman_all_2004}.

Let $X: \Omega \rightarrow \R$ denote a continuous random variable mapping a sample space of events $\Omega$ to the real numbers. The probability density function (pdf) of $X$ is denoted by $f_X(x)$. In general, we will drop the random variable and write the pdf as $f(x)$. 

The pdf $f(x)$ satisfies the properties:
\begin{enumerate}
\item $\int_{-\infty}^{\infty}f(x)dx =1$,
\item $P(a < X < b) = \int_{a}^{b}f(x)dx$,
\end{enumerate}
for $a\leq b$. Note that $P(a < X < b)$ is the probability that the random variable $X\in(a, b)$.

The corresponding cumulative distribution function (cdf) $F_{X}(x)$ is
\begin{equation}
F_{X}(x) = \int_{-\infty}^{x}f_{X}(x)dx \equiv P(X \leq x).
\end{equation}
Again, we write the cdf as $F(x)$.

The expectation of a random variable is
\begin{equation}
    E(X) = \int_{-\infty}^{\infty}xf(x)dx.
    \label{eq:expectation}
\end{equation}
The notation $X\sim f(x)$ is used to indicate that the random variable $X$ has distribution $f$. By~\eqref{eq:expectation}, we have that 

\begin{equation}
    E(g(X)) = \int_{-\infty}^{\infty}g(x)f(x)dx.
\end{equation}

Let $p(x)$ denote a distribution and let $X\sim p(x)$. We introduce the notation 
\begin{equation}
    E_{X\sim p(x)}(f(X)) = \int_{-\infty}^{\infty}p(x)f(x)dx.
\end{equation}
In the multivariate case, i.e., when $X=(X_1, X_{2},\ldots, X_n)^{\top}$ and each $X_{i}$ is a random variable, the pdf $f_{X}(x)$ is a joint distribution. The marginal distribution for $X_{i}$ is then defined as
\begin{equation}
    f_{X_{i}}(x_{i}) = \int_{-\infty}^{\infty}f_{X}(x)\prod_{\substack{j=1 \\j\neq i}}^{n} dx_{j}.
\end{equation}
For independent and identically distributed (IID) random variables we have the property that the joint distribution factors as the product of the univariate distributions, i.e.,
\begin{equation}
    f_{X}(x) = \prod_{i=1}^{n}f_{X_{i}}(x_{i}).
\end{equation}
Consequently, 
\begin{equation}
    \log{f_{X}(x)} = \sum_{i=1}^{n}\log{f_{X_{i}}(x_{i})}.
\end{equation}

Given two random variables $X$ and $Y$, where $f_{Y}(y)>0$, the conditional distribution of $X$ given $Y$ is
\begin{equation}
    f_{X\mid Y}(x\mid y) = \frac{f_{X,Y}(x,y)}{f_{Y}(y)}.
    \label{eq:conditionalpdf}
\end{equation}
This shows that the conditional pdf $f_{X\mid Y}(x\mid y)$ is the ratio of the joint distribution over the marginal distribution of $Y$.

An immediate consequence of~\eqref{eq:conditionalpdf} is  Bayes' theorem.
\begin{theorem}
For two continuous random variables $X$ and $Y$
\begin{equation}
    f(x\mid y) = \frac{f(y\mid x)f(x)}{f(y)},
    \label{eq:Bayes}
\end{equation}
provided $f(y) > 0$.
\end{theorem}

For two distributions $p(x)$ and $q(x)$, we define the Kullback-Leibler divergence

\begin{equation}
    D_{KL}(p \mid\mid q) = \int_{-\infty}^{\infty}p(x)\log{\frac{p(x)}{q(x)}}.
    \label{eq:KLDivergence}
\end{equation}

A parametric distribution $p(x;\theta)$ is a set (or family) of distributions that are fully specified by a finite set of parameters $\theta$. An example is the normal distribution $N(\mu, \sigma^{2})$ which is a family of distributions defined by the mean $\mu$ and variance $\sigma^{2}$, i.e., $f(x)\sim N(\mu, \sigma^{2})$ provided

\begin{equation}
    f(x) = \frac{1}{\sigma\sqrt{2\pi}}\int_{-\infty}^{\infty}e^{\frac{-(x-\mu)^{2}}{\sigma^{2}}}.
\end{equation}

Generic parameterized distributions are also denoted as $p_{\theta}(x)$.

Let $X$ and $Y$ be (continuous) random variables with joint probability distribution $f(x,y)$. The mutual information of $X$ and $Y$ is defined as

\begin{equation}
I(X,Y) = \int\int f(x,y)\log\left(\frac{f(x,y)}{f(x)f(y)}\right).
\label{eq:MutualInformation}
\end{equation}

For discrete random variables $X$ and $Y$ the formula for mutual information becomes

\begin{equation}
I(X,Y) = \sum_{x}\sum_{y} f(x,y)\log\left(\frac{f(x,y)}{f(x)f(y)}\right).
\label{eq:MutualInformation}
\end{equation}

Let $X$ be a continuous random variable with probability distribution $f(x)$, then the differential entropy $H(X)$ is defined as

\begin{equation}
H(X) = -\int f(x)\log(f(x)) dx.
\label{eq:DifferentialEntropy}
\end{equation}

For discrete random variables, we obtain the Shannon entropy

\begin{equation}
H(X) = -\sum_{x}f(x)\log(f(x)).
\label{eq:ShannonEntropy}
\end{equation}


\section{Variational Autoencoders}
Variational autoencoders (VAEs) introduced in \cite{kingma_auto-encoding_2014} are generative models that consist of an encoder and a decoder. The encoder maps input data to a latent space that is characterized by a distribution. A sampled output of this distribution is then fed into a decoder, which generates the output of the model.

To make this precise, we consider a dataset $\{x^{(1)},\ldots,x^{(i)},\ldots, x^{(N)}, \}$ consisting of $N$ independently and identically distributed (i.i.d) samples of a continuous random variable $X$. This data is generated by a random process that involves an unobserved continuous random variable $Z$. This is a two step process:
\begin{enumerate}
    \item a value $z^{(i)}$ is generated from a prior distribution $p_{\theta^{*}}(z)$, and
    \item a value $x^{(i)}$ is generated from a conditional distribution $p_{\theta^{*}}(x\mid z)$.
\end{enumerate}
The prior and conditional distributions come from parametric families of distributions $p_{\theta}(z)$ and $p_{\theta}(x\mid z)$, respectively.

Since we do not have access to the true posterior, $p_{\theta}(z\mid x)$, we approximate it using a recognition model (encoder) $q_{\phi}(z\mid x)$. For VAEs, the encoder $q_{\phi}(z\mid x)$ and decoder $p_{\theta}(x\mid z)$ models are represented by deep neural networks. We outline the procedure for determining the parameters $(\phi, \theta)$ of these models via training.

For a set of IID, the log of the marginal likelihood decomposes as the sum of the logs of the individual marginal likelihoods, i.e.,

\begin{equation}
    \log(p_{\theta}(x^{(1)},\ldots, x^{(n)})) = \sum_{i=1}^{n}\log(p_{\theta}(x^{(i)})).
\end{equation}

We now introduce the following proposition

\begin{proposition}
    The marginal probability $\log(p_{\theta}(x))$ is equivalently expressed as
    \begin{align}
        \log(p_{\theta}(x)) &= E_{q_{\phi}(z\mid x)}\left(\log\left(\frac{p_{\theta}(x,z)}{q_{\phi}(z\mid x)}\right) \right) + D_{KL}(q_{\phi}(z\mid x) \mid\mid p_{\theta}(z\mid x)) \\
        &\equiv L_{\theta,\phi}(x) +  D_{KL}(q_{\phi}(z\mid x) \mid\mid p_{\theta}(z\mid x)). 
    \end{align}
The term $ L_{\theta,\phi}(x) = E_{q_{\phi}(z\mid x)}\left(\log\left(\frac{p_{\theta}(x,z)}{q_{\phi}(z\mid x)}\right) \right)$ is the evidence lower bound and $D_{KL}(q_{\phi}(z\mid x) \mid\mid p_{\theta}(z\mid x))$ is the Kullback-Leibler~\eqref{eq:KLDivergence} divergence between the encoder and decoder models.
\end{proposition}

\begin{proof}
    We have that

    \begin{align}
        \log(p_{\theta}(x)) &= \int_{z}q_{\phi}(z\mid x)\log(p_{\theta}(x))dz \\
        &= E_{q_{\phi}(z\mid x)}\left(\log(p_{\theta}(x))\right) \\
        &= E_{q_{\phi}(z\mid x)}\left(\log\left(\frac{p_{\theta}(x,z)}{p_{\theta}(z\mid x)}\right)\right)\label{eq:contbayes} \\
        &= E_{q_{\phi}(z\mid x)}\left(\log\left(\frac{p_{\theta}(x,z)}{q_{\phi}(z\mid x)}\frac{q_{\phi}(z\mid x)}{p_{\theta}(z\mid x)}\right)\right) \\
        &= E_{q_{\phi}(z\mid x)}\left(\log\left(\frac{p_{\theta}(x,z)}{q_{\phi}(z\mid x)}\right) \right) + E_{q_{\phi}(z\mid x)}\left(\log\left(\frac{q_{\phi}(z\mid x)}{p_{\theta}(z\mid x)}\right) \right) \label{eq:linearity}
    \end{align}

    Observe that \eqref{eq:contbayes} follows from Bayes theorem and \eqref{eq:linearity} follows by properties of logs and the linearity of the expectation. The first term in \eqref{eq:linearity} is the Evidence Lower Bound (ELBO). The second term is the non-negative KL divergence between $q_{\phi}(z\mid x)$ and $p_{\theta}(z\mid x)$.
\end{proof}

The non-negativity of the KL divergence implies that $L_{\phi,\theta}(x)\leq \log(p_{\theta}(x))$. In VAEs, we try to maximize the ELBO, which requires differentiating the ELBO with respect to both $\theta$ and $\phi$.


\section{Text Style Transfer}
The goal of text style transfer is to transform a piece of text in one particular style into a different style, while preserving the meaning of the original text. More precisely, given an input text $x$ of style $s$ and a corresponding target style $s^{\prime}$, generate text $x^{\prime}$ for which the original content in $x$ is maintained in $x^{\prime}$.

We then model the generation of $x^{\prime}$ as a conditional probability, i.e.,

\begin{equation}
p(x^{\prime}\mid x,s,s^{\prime})
\end{equation}



\subsection{Style as a Latent Variable}


\bibliography{StyleTransfer.bib}
\end{document}
